\scanpage{32}

\figure{\begin{tikzpicture}[x=1cm,y=1cm,every node/.style={anchor=north west,align=left}]
\draw (0.35,0.15)--(0,0.15)--(0,-12.55) to[out=-90,in=180] (0.3,-12.8);
\draw (0,-11.15)--(0.4,-11.15);
\node[text width=13.6cm] at (0.4,0) {1. Généralisation des th. de finitude formelle, et du th. de comparaison. Application : un th. d'existence};
\node[text width=13.6cm] at (0.4,-1.35) {2. Le Théorème de finitude locale et ses applications.};
\node[text width=13.6cm] at (0.4,-2) {\deleted{3. Fermés de Lefschetz-Grauert; Remarques.}};
\node[text width=13.6cm] at (0.4,-2.9) {\deleted{4. \unsure{Premières} propriétés des \uline{Exemples} \unsure{fondamentaux}}};
\node at (10.5,-2.65) {\inserted{1}};
\draw[rounded corners] (0.35,-1.98) rectangle (12.4,-3.42);
\draw (0.45,-3.3)--(1.25,-2.1)--(1.9,-3.35)--(2.55,-2.1)--(3.35,-3.35)--(4,-2.1)--(4.7,-3.35)--(5.4,-2.1)--(6.3,-3.35)--(7,-2.1)--(7.9,-3.35)--(8.7,-2.1)--(9.5,-3.35)--(10.3,-2.1)--(11.1,-3.35)--(12.1,-2.2);
\node[text width=13.6cm] at (0.65,-3.85) {3. Cas particulier \unsure{I|II} : \deleted{Étude}\inserted{\unsure{théorie}} dans \deleted{\unsure{les}} schémas locaux.};
\node[text width=13.6cm] at (0.65,-4.85) {4. Cas particulier 2 : \deleted{Étude}\inserted{\unsure{théorie}} des schémas projectifs \inserted{\unsure{quelconques|réguliers}}};
\node[text width=13.6cm] at (0.65,-6.15) {5. Fermés de Lefschetz-Grauert. Remarques, exemples.};
\node at (0.65,-7.2) {6. Propriétés liées aux $\pi_0,$};
\node (pi) at (7.1,-7.2) {$\pi_1$};
\node at (7.7,-7.2) {, Pic.};
\draw (7.45,-7.48) ellipse[x radius=.32,y radius=.5];
\node at (0.35,-8.25) {[7. Propriétés de permanence]};
\node (eight) at (0.65,-9.5) {8.};
\draw (0.9,-9.7) ellipse[x radius=.46,y radius=.55];
\node[text width=12cm] at (1.65,-9.45) {Théorèmes de pureté, et Théorèmes de Lefschetz pour le $\pi_1$.};
\node[draw,rounded corners,rotate=28,text width=3.1cm] (m) at (-2,-6.65) {\margin{\unsure{utilise|utiliser} IV, \S\ sur les morphismes étales\ldots}};
\draw[->] (m.south east) to[out=12,in=175] (7.25,-6.95);
\draw[->] (-.95,-7.9) to[out=-55,in=150] (0.45,-9.45);
\node at (0.7,-11.45) {\deleted{\unsure{9}}};
\draw (1,-11.8) ellipse[x radius=.45,y radius=.5];
\draw (0.65,-11.45)--(1.3,-11.8)--(0.65,-12)--(1.25,-11.5)--(.75,-12.15)--(1.3,-11.65);
\node at (0.7,-12.35) {9};
\node[text width=12cm] at (2,-12) {\uline{Anneaux} \deleted{\illegible{1}} \uline{quasi-factoriels} et \uline{th de} \\ \uline{Lefschetz pour Pic.}};
\end{tikzpicture}}

\margin{31}
